\documentclass[a4paper, czech]{article}

\title{
    \begin{framed}
        \setstretch{1.5}
        {\large Audio elektronika (BPC/BKC-AUD)}\\
        Laboratorní úloha č. 3 (protokol)\\
        \textbf{Koncový audio zesilovač ve třídě T}
    \end{framed}
    }
\author{Karolína A. Šebestová}
\date{3. 3. 2026, 11:00}

\usepackage[czech]{babel}
\usepackage{indentfirst}
\usepackage{graphicx}
\usepackage{float}
\usepackage[margin=1.5cm]{geometry}
\usepackage{booktabs}
\usepackage{amsmath}
\usepackage{multirow}
\usepackage{colortbl}
\usepackage{framed} % \begin{framed}
\usepackage{setspace} % \setstretch{}
\usepackage{siunitx}
\usepackage{multirow}
\usepackage{xcolor}

\begin{document}

\maketitle

\section{Měření}

\subsection{Měření modulové kmitočtové charakteristiky}

\begin{equation}
    A_\text{U} = \textcolor{teal}{20 \cdot \log \frac{U_2}{U_1}} = 20 \cdot \log \frac{0,98}{100 \cdot 10^{-3}} = \underline{\underline{19,82\,\text{dB}}}
\end{equation}

\begin{table}[H]
    \catcode`\-=12
    \centering
    \caption{Modulová kmitočtová charakteristika ($U_1$ = 100\,mV).}

    \begin{tabular}{ccc}
        \toprule
        $f$ [Hz] & $U_2$ [mV] & $A_\text{U}$ [dB] \\
        \cmidrule(rl){1-3}
        10     & 0,98  & 19,82  \\
        20     & 1,02  & 20,17  \\
        30     & 1,03  & 20,26  \\
        50     & 1,03  & 20,26  \\
        70     & 1,03  & 20,26  \\
        100    & 1,04  & 20,34  \\
        200    & 1,04  & 20,34  \\
        300    & 1,04  & 20,34  \\
        500    & 1,04  & 20,34  \\
        700    & 1,04  & 20,34  \\
        1000   & 1,04  & 20,34  \\
        2000   & 1,04  & 20,34  \\
        3000   & 1,04  & 20,34  \\
        5000   & 1,05  & 20,42  \\
        7000   & 1,07  & 20,59  \\
        10000  & 1,09  & 20,75  \\
        20000  & 1,23  & 21,80  \\
        30000  & 1,37  & 22,73  \\
        50000  & 1,28  & 22,14  \\
        100000 & 0,016 & -15,92 \\
        \toprule
    \end{tabular}
\end{table}

\begin{figure}[H]
    \centering
    \includegraphics[width=\textwidth]{grafy/graf1.pdf}
    \caption{Modulová kmitočtová charakteristika koncového zesilovače ve třídě T.}
\end{figure}

Přibližná šířka pásma zesilovače $B = 52490\,\text{Hz}$.

\subsection{Měření maximálního výstupního výkonu pro limitaci}

$U_{\text{výst}} = 5,81\,\text{V}$

\begin{equation}
    P_{\text{max}} = \textcolor{teal}{\frac{U_{\text{výst}}^2}{R_\text{Z}}} = \frac{5,81^2}{8} = \underline{\underline{4,22\,\text{W}}}
\end{equation}

Maximální výkon $P_{\text{max}} = 4,22\,\text{W}$.

\subsection{Měření účinnosti}

$U_\text{N} = 11,67\,\text{V}$

$I_\text{N} = 0,52\,\text{A}$

\begin{equation}
    \eta = \textcolor{teal}{\frac{P_\text{max}}{U_\text{N} \cdot I_\text{N}}} = \frac{4,22}{11,67 \cdot 0,52} = \underline{\underline{0,6953}}
\end{equation}

Účinnost $\eta = 69,53\,\%$.

\subsection{Měření vstupního odporu zesilovače}

Vstupní odpor $R_\text{vst} = 22\,k\Omega$

\subsection{Měření výstupního odporu zesilovače}

$R_\text{Z} = 8\,\Omega$

$U_2 = 1,08\,\text{V}$

$U_{20} = 1,19\,\text{V}$

\begin{equation}
    R_\text{výst} = \textcolor{teal}{R_\text{Z} \frac{U_{20} - U_2}{U_2}} = 8 \cdot \frac{1,19 - 1,08}{1,08} = \underline{\underline{0,81\,\Omega}}
\end{equation}

Výstupní odpor $R_\text{výst} = 0,81\,\Omega$.

\subsection{Určení činitele tlumení zátěže}

\begin{equation}
    D = \textcolor{teal}{\frac{R_\text{Z}}{R_\text{výst}}} = \frac{8}{0,81} = \underline{\underline{9,82}}
\end{equation}

Činitel tlumení zátěže $D = 9,82$.

\subsection{Měření rychlosti přeběhu}

$U_1 = 0,535\,\text{V}$

Rychlost přeběhu náběžné hrany $SR_\text{rise} = 1,92\,\text{V/}\mu\text{s}\ (\Delta U = 10,37\,\text{V},\ \Delta t = 5,4\,\mu\text{s}).$

Rychlost přeběhu sestupné hrany $SR_\text{fall} = 2,16\,\text{V/}\mu\text{s}\ (\Delta U = 10,37\,\text{V},\ \Delta t = 4,8\,\mu\text{s}).$


\subsection{Měření harmonického zkreslení}

\begin{table}[H]
    \catcode`\-=12
    \centering
    \caption{THD koncového zesilovače ve třídě T ($f$ = 1\,kHz, $R_\text{Z}$ = 8\,$\Omega$, $THD+N_\text{GEN}$ = 0,023\%).}
    \begin{tabular}{ccccc}
        \toprule
        $U_2$ [V] & $P_2$ [W] & $k_2$ [\%] & $k_3$ [\%] & $THD+N$ [\%] \\
        \cmidrule(rl){1-5}
        0,55   & 0,0378125 & 0,0163  & 0,0391  & 0,15       \\
        2,75   & 0,9453125 & 0,2034  & 0,015   & 0,275      \\
        4,96   & 3,0752    & 0,0155  & 0,0421  & 0,219      \\
        5,85   & 4,2778125 & 0,0202  & 0,0517  & 0,104      \\
        7,19   & 6,4620125 & 0,1186  & 0,2519  & 1,389      \\
        \bottomrule
    \end{tabular}
\end{table}

\begin{figure}[H]
    \centering
    \includegraphics[width=\textwidth]{grafy/graf2.pdf}
    \caption{Měření harmonického zkreslení na výstupním výkonu.}
\end{figure}


\section{Použité měřící přístroje}

\begin{itemize}
    \item \textbf{GEN} nízkofrekvenční funkční generátor Agilent 33220A
    \item \textbf{NMV} nízkofrekvenční milivoltmetr Grundig MV100
    \item \textbf{OSC} digitální osciloskop Agilent 54621A
    \item \textbf{NZ} napájecí zdroj MCP M10-DP-305E (napájení přípravku +15V)
    \item \textbf{A} ampérmetr DMM-3900
    \item \textbf{V} voltmetr DMM-3900
    \item měřený přípravek „Koncový nízkofrekvenční zesilovač ve třídě T“
    \item vstupní cejchovaný proměnný odpor 100\,k$\Omega$ a propojovací vodiče 4 x BNC – „banánek“
\end{itemize}

\section{Závěr}

Dolní pokles o 3\,dB nelze z naměřené modulové kmitočtové charakteristiky odečíst.
Odečetli jsme ho proto odhadem.
Horní pokles o 3\,dB se nachází v oblasti 50 kHz a 100 kHz, získání přesnější hodnoty by vyžadovalo větší množství měření v dané oblasti.
Zesilovač v propustné oblasti zesiluje přibližně kolem 21\,dB.

Maximální výstupní výkon zesilovače v lineární oblasti je 4,22\,W.
Z této hodnoty jsme vypočetli účinnost 69,53\,\%.

Vstupní odpor zesilovače při kmitočtu 1\,kHz je 22\,k$\Omega$.
Poté byl vypočítán i výstupní odpor zesilovače 0,81\,$\Omega$.
Zesilovač má velký vstupní a malý výstupní odpor, pracuje proto tedy jako zdroj napětí.
Činitel tlumení zátěže je 9,82 a odpovídá tedy teoretickým předpokladům.

Zesilovač reaguje na náběžnou hranu s rychlostí 1,92\,$\text{V/}\mu\text{s}$ a na sestupnou s rychlostí 2,16\,$\text{V/}\mu\text{s}$.
Měření bylo pro vyšší přesnost provedeno manuálně pomocí kurzorů na osciloskopu.

Harmonické zkreslení $THD+N$ se začíná projevovat po hodnotě výkonu 5\,W.
Hodnoty $k_2$ a $k_3$ se však drží pod hodnotou 1\,\% i pro vyšší výstupní výkon.
Zkreslení vlastního generátoru splňuje podmínku pro poměr k výsledné hodnotě $THD+N$.

\end{document}